\section{A historical context of electromagnetic optics and polarisation}

The idea that light is an electromagnetic phenomenon --- and that its transverse vibration
carries a hidden orientation we call polarisation --- is one of the great unifications of
nineteenth-century physics. Before light could be understood as a travelling disturbance of
electric and magnetic fields, it had to shed two older skins: Newton's corpuscles and the
purely mechanical aether of the early wave theorists. The story of polarisation, in
particular, is the story of how a seemingly minor curiosity --- light reflected from a
window behaving differently from ordinary light --- forced physicists to accept that the
optical vibration is transverse, and ultimately vectorial.

The first quantitative encounter with polarisation came in 1808, when Étienne-Louis Malus,
looking through a calcite crystal at sunlight reflected from the windows of the Luxembourg
Palace in Paris, noticed that the intensity of the transmitted beam changed as he rotated
the crystal. From these observations he extracted the law that bears his name, $I =
I_0\cos^2\theta$, relating the transmitted intensity to the angle between the analyser and
the plane of polarisation \cite{malus1809theorie}. Malus had discovered that reflection
could endow light with the same directional property that double refraction revealed,
without invoking any crystal at all. A few years later, in 1815, David Brewster
established the law that the polarisation upon reflection is complete at a particular angle,
the Brewster angle, for which the reflected and refracted rays are perpendicular
\cite{brewster1815laws}.

These facts were deeply puzzling within the longitudinal wave picture inherited from sound.
The resolution came from Augustin-Jean Fresnel and François Arago, who between 1816 and 1819
showed experimentally that two beams polarised at right angles do not interfere. Fresnel
drew the inevitable, audacious conclusion: the optical vibration must be \emph{transverse}
to the direction of propagation. From this single hypothesis he derived the laws governing
the amplitudes of reflected and refracted light --- the Fresnel equations --- which remain
in use today. Transversality made polarisation a natural attribute of light: a wave
oscillating in a plane perpendicular to its motion can do so along infinitely many
directions, and the choice of that direction \emph{is} its state of polarisation.

A decisive step toward a quantitative description of partially polarised light was taken by
George Gabriel Stokes in 1852. Recognising that realistic beams are never perfectly
polarised, Stokes introduced four measurable quantities --- today called the Stokes
parameters --- that completely characterise the state of polarisation, including the degree
of polarisation, of any beam, coherent or not \cite{stokes1852composition}. Decades later,
Henri Poincaré gave these parameters a beautiful geometric home: every state of
polarisation corresponds to a point on the surface of a sphere, the Poincaré sphere, on
which the poles represent circular polarisation and the equator linear polarisation. This
geometric picture is the one that reappears, in the scalar reduction of diffraction theory,
when one decomposes the field onto two opposite states of the sphere.

The grand synthesis arrived with James Clerk Maxwell. In his 1865 memoir \textit{A
Dynamical Theory of the Electromagnetic Field} \cite{maxwell1865dynamical}, Maxwell showed
that the equations governing electricity and magnetism admit wave solutions propagating at
a speed numerically equal to the measured speed of light, and concluded that ``light is an
electromagnetic disturbance.'' Polarisation was thereby identified with the direction of
the electric field vector $\mathbf{E}$, and the whole edifice of physical optics became a
chapter of electromagnetism. The experimental confirmation came in 1887, when Heinrich
Hertz generated and detected electromagnetic waves in the laboratory, verifying that they
reflect, refract and polarise exactly as light does \cite{hertz1893electric}.

In the twentieth century the description of polarisation was recast in the language of
linear algebra. R.\ Clark Jones, in a celebrated series of papers beginning in 1941,
introduced the calculus of two-component complex vectors and $2\times2$ matrices that now
bears his name, ideally suited to fully polarised, coherent light \cite{jones1941new}. For
partially polarised or depolarising systems, the complementary $4\times4$ Mueller calculus
acting on the Stokes vector provides the appropriate tool. The modern treatment of all
these matters --- transversality, the Fresnel equations, the Stokes--Poincaré formalism and
the Jones and Mueller calculi --- is presented with classic completeness in Born and Wolf's
\textit{Principles of Optics} \cite{born1999principles}, the reference against which later
treatments are still measured.

From Malus's rotating crystal to the polarisation-encoded qubits of quantum information,
the orientation of the optical field has proven to be far more than an ornament: it is a
genuine degree of freedom of light, vectorial in nature, and the bridge between the scalar
optics of rays and the full electromagnetic description developed in the chapters that
follow.
